\selectlanguage{english}
\phantomsection\addcontentsline{toc}{section}{Overview of these pages}
\nocite{GrothendieckDerivateursVI}
\noindent\textbf{Integration, local constancy, and the choice of weak equivalences.}

The last part of Chapter VI turns the construction of relative homotopy types into a calculus. Its first question is concrete. A category over $S'$ can also be regarded as a category over $S$ by composing its structural functor with $f:S'\to S$. This elementary operation already gives the homological direct image $f_!$ after localization. The difficulty lies on the other side of the proposed adjunction: ordinary pullback does not preserve the weak equivalences used to define relative homotopy types. Grothendieck therefore represents inverse image by a replacement followed by ordinary pullback. The argument on manuscript pages 138--146 identifies the unit and counit of the resulting adjunction and checks their two triangular identities.

It helps to fix the variance before interpreting these formulas. For $p:X\to S$, the category $s\backslash X=(s\downarrow p)$ has objects $(x,a:s\to p(x))$. The replacement $\Phi_SX$ is the Grothendieck construction of the contravariant diagram $s\mapsto s\backslash X$. Write $Q_SX$ for the image of $X$ in
\[
 \HOT_W(S)=(\Cat/S)[(W_S^{\mathrm d})^{-1}],\qquad
 W_S^{\mathrm d}=\{u\mid s\backslash u\in W\text{ for every }s\in S\}.
\]
Then $f_!Q_{S'}X'=Q_SX'$, whereas
\[
 f^*Q_SX=Q_{S'}(\Phi_SX\times_SS').
\]
The replacement map $i_X:X\to\Phi_SX$ belongs to $W_S^{\mathrm d}$. The counit is represented by the composite $(i_X)^{-1}\pr_1$, with the inverse taken \emph{after localization}. Attempting to invert the pullback of $i_X$ would require precisely the invariance that has not been established. The abandoned attempt on page 142 and its repair on page 143 show why the order of these operations matters.

The next difficulty concerns products. If both representatives are smooth over $S$, their ordinary fibre product represents the product of their localized classes. One smooth representative does not suffice when the other is arbitrary. Grothendieck records the false extension, finds a counterexample over $[1]=(0\to1)$, and replaces it by a statement about an endofunctor. For every fixed $X\to S$, the construction
\[
 M_{X/S}(Q_SY)=Q_S(X\times_S\Phi_SY)
\]
is defined, although it need not depend only on the localized class $Q_SX$. The universal formula is $f_!f^*=M_{S'/S}$. Identifying its right side with multiplication by $Q_SS'$ requires an additional hypothesis. This distinction returns on pages 199--202, where Grothendieck asks whether a weak $W$-fibration factor suffices even when the other raw representative is arbitrary. We prove that it does. If a split replacement $E$ is classified by a locally constant diagram $F:S^{\mathrm{op}}\to\Cat$, transport gives natural $W$-equivalences
\[
s\backslash(E\times_S Y)\longrightarrow F(s)\times(s\backslash Y).
\]
They show that multiplication by $E$ preserves colocal equivalences viewed over $S$. Reduction from $X$ to $E$ completes the product formula under the basic localizer hypotheses. The stronger assertion of colocal equivalence viewed over $E$ is unnecessary; that is the distinction which the unfinished argument misses.

The base-change theorem on pages 162--169 is the geometric culmination of this first movement. In a cartesian square
\[
\begin{tikzcd}
 X'\arrow[r,"q"]\arrow[d,"f'"']&X\arrow[d,"f"]\\
 S'\arrow[r,"p"']&S,
\end{tikzcd}
\]
the canonical mate $f'_!q^*\to p^*f_!$ is invertible when $f$ is smooth or $p$ is proper, with the corresponding relative versions for $W$. Smoothness allows the same representative to compute both sides. Properness ensures that strict pullback preserves the colocal weak equivalences. Grothendieck first tries to obtain the latter assertion from a derivator theorem, then notices that the argument would use the unproved conservativity axiom Der~2. His direct proof on pages 168--169 removes that dependence: a suitable map of comma categories is a universal $W$-equivalence, and a commutative square with its vertical arrows in $W$ transfers membership in $W$ between the horizontal arrows. Thus the acknowledged gap in the detour does not invalidate the final base-change proof.

This continues the discussion of mates and exact squares indexed in DER\_I-a, at manuscript pages 31 and 47--58. Here the issue is the construction of the particular theory $\HOT_W$. The general Der~2 gap was later repaired: comparison with diagrams of simplicial sets and joint conservativity of their point evaluations apply to every fundamental localizer, without an accessibility assumption \citep[Propositions~4.2.7 and~3.1.30]{Cisinski2006}. The commentary states this comparison explicitly, while distinguishing it from the proof actually present. Modern homotopy Kan extensions provide the surrounding framework \citep[\S1]{Groth2013}.

\medskip\noindent\textbf{The revision beginning on page 171.}

Grothendieck now asks how much of the theory depends on ordinary weak homotopy equivalences and how much uses only the axioms on $W$. The notation becomes more demanding because several inequivalent tests must be kept visible. A map over $S$ may be a weak equivalence of total categories, a weak equivalence on every strict fibre, or a weak equivalence after passing to the over- or under-comma categories. Smooth objects identify the fibrewise and colocal tests; proper objects identify the fibrewise and local tests. Weak $W$-fibrations identify both local tests with the fibrewise test. The reverse implication from a total equivalence to a fibrewise equivalence is an additional axiom, W(4b), not a consequence of those identifications.

Accordingly the revised account has four canonical localized categories. Besides $\HOT_W(S)$ and its opposite-variance companion, it distinguishes
\[
 \HOT_W^{\mathrm{lc}}(S)=\Fib_WS[(W_S^{\mathrm f})^{-1}]
 \quad\hbox{and}\quad
 \HOT_W^{\mathrm{lc}!}(S)=\Fib_WS[W_S^{-1}].
\]
There is a comparison $\lambda$ from the first to the second. W(4b) makes their defining weak equivalences identical and hence makes $\lambda$ an equivalence. Without that axiom, a map of total categories can erase fibre information. The worked homological example in the commentary makes this loss explicit; it also explains why an equivalence after a further localization cannot by itself justify arbitrary base change before that localization.

A second question is which representatives suffice. The replacements $\Phi$ and $\Psi$ supply smooth and proper representatives in their respective variances. The path construction $C$ supplies perfect representatives, but its comparison map is initially only a total weak equivalence. Improving the representative and controlling the weak equivalence used to improve it are different assertions. Much of pages 176--186 is devoted to preventing the second assertion from being silently inferred from the first.

The modern answer settles several historical speculations. Cisinski proved that ordinary weak homotopy equivalences form even the \emph{weak} fundamental localizer of smallest size \citep[Theorem~6.1.18]{Cisinski2006}; the corresponding theorem for fundamental localizers appeared earlier \citep[Theorem~2.2.11]{Cisinski2004Minimal}. Strong saturation and stability under arbitrary filtered colimits are automatic for fundamental localizers \citep[Proposition~4.2.4 and Corollary~4.2.22]{Cisinski2006}. Thus the three proposed minima $W_\omega$, $W'_\omega$, $W''_\omega$ are all $W_\infty$. Together with Chapter~XVI, Theorem~4.4.13, these results also identify the original perfect-object test W(4b$'$) with W(4b) under the revised hypotheses. More strongly, W(4b) forces a nontrivial fundamental localizer to be $W_\infty$ \citep[Theorem~6.4.22 and Remark~6.4.23]{Cisinski2006}. This does not collapse the four-category construction for arbitrary $W$: rational homology equivalences form a larger localizer and fail the fibre axiom, as the two worked examples demonstrate.

\medskip\noindent\textbf{Perfection by iteration and the questions left by the text.}

The failure of a single replacement is especially sharp over the branching poset $S=(0\to a,0\to b)$. The comma replacement $\Phi_SX$ can be proper over this base only when $X$ is empty, for a nondegenerate localizer. We spell out the empty-comma-category obstruction, including the case of objects lying over $0$. This does not prevent simultaneous smooth and proper replacement. It explains why the operation that produces it must be more substantial than one application of $\Phi$.

The proposal on pages 192--195 alternates the two replacements:
\[
 X\longrightarrow\Phi X\longrightarrow\Psi\Phi X
 \longrightarrow\Phi\Psi\Phi X\longrightarrow\cdots.
\]
Under W(6 bis), which here concerns sequential colimits, the colimit is smooth by its odd subsequence and proper by its even subsequence. For a weak $W$-fibration, every comparison in the sequence is fibrewise a weak equivalence, and so is the comparison with the colimit. This yields the desired equivalence after fine localization between weak fibrations and perfect representatives. Modern filtered-colimit closure makes the extra hypothesis automatic. Since the minimum classes coincide, every full intermediate category containing $\Parf_\omega S$ gives the same fine localization; the speculative proliferation of inequivalent fine presentations is therefore resolved. The fine--coarse comparison $\lambda$ remains a separate question. Nor does the construction become a literal Grothendieck bifibration: cartesian and cocartesian lifting properties are stronger categorical conditions. We prove essential surjectivity of strict bifibration representatives over bases having an initial or terminal object. For the unrestricted literal question, we report the limit of the checked literature, without declaring a globally open problem.

The manuscript later declines to check a triangular identity needed for the reflector obtained from this iteration. We complete the argument by an elementary localization lemma: for a coaugmented endofunctor $(T,\eta)$, if $T\eta$ and $\eta T$ are invertible, then precomposition with $\eta_X$ identifies maps from $TX$ to every $\eta$-local object with maps from $X$ to that object. This both constructs the reflection and forces the two comparison maps to agree. The additional results marked \ours\ are explicit continuations of arguments in the supplied pages; the mark asserts responsibility for the proof given here, not historical priority.

Grothendieck's own later treatment in Chapter XVI, especially Lemma~4.4.9, Corollaries~4.4.10--4.4.12 and Theorem~4.4.13, is indispensable for reading the retrospective annotations \citep{GrothendieckDerivateursXVI}. The development of locally constant diagrams in model categories and derivators supplies a further precise comparison \citep[Theorem~1.3; \S\S2--3]{Cisinski2009LocallyConstant}. In the ordinary space-valued case, locally constant coherent diagrams are governed by the groupoid completion of the indexing category in the homotopical sense; replacing that object by its fundamental \emph{one}-groupoid would discard higher information. The audit does not identify this modern description with a strict bifibration theorem without an additional argument.

The surviving text ends during a restatement of base change. Manuscript page 170 is absent, and the final words explicitly announce that the continuation is missing. We preserve both boundaries. The audit records a statement as proved, sketched, conjectured, or withdrawn according to the argument actually present, and places any later repair in a separate field. In particular Der~2 is not credited to Theorem~3 merely because its name remains in the printed theorem: the attached note expressly retracts that part of the claim.
